\section{The compact graph construction}\label{sec:compact-construction}

Fix the compact data and the asymptotically flat attachment supplied by
Proposition~\ref{prop:af-attachment}. Thus the connected compact inner
manifold $(\Omega_i,g)$ has boundary $S\sqcup\mathcal C$, where
$\mathcal C$ is a sphere, and carries a smooth negative definite tensor
$K$. Put $\tau=\tr_g K$. The smooth exterior $(\Omega_o,g_o)$ is scalar
flat and asymptotically flat. The induced metrics on $\mathcal C$ agree.
All normals $n$ at $S$ and $\mathcal C$ point towards the designated end;
in particular the outward normal of $\Omega_i$ at $S$ is $-n$. Write
$k_T=\tr_{T\mathcal C}K$ on the seam and use the analogous notation on
$S$. The fixed data satisfy
\begin{equation}\label{eq:compact-base-data}
 \begin{gathered}
 H_g(S)\le0,\qquad H_i>|k_T|,\qquad
 H_o=\sqrt{H_i^2-k_T^2},\\
 D(v):=R_g+\tau^2-|K|^2-2(\divg_gK-\dd\tau)(v)>0
 \quad (|v|_g\le1).
 \end{gathered}
\end{equation}
No deformation parameter changes these base data. In this section every
integral without a metric subscript uses the base metric on its own side.
Piecewise integrals mean the sum of the two such integrals. Constants may
depend on these fixed smooth data.

On $\Omega_i$ define the fields, for $N_i>0$, by
\begin{equation}\label{eq:compact-constitutive}
 \begin{gathered}
 p=s^2\nabla f,\qquad s^2+|\dd f|^2s^4=N_i^2,\qquad v=p/N_i,\\
 z=|v|,\quad y=z^2,\quad d=\sqrt{1-y},\quad W=d^{-1}=N_i/s.
 \end{gathered}
\end{equation}
The positive root specifies $s$ smoothly, including at $\dd f=0$.
On a constant-height boundary we write $\widehat z=v\cdot n$.
On $\Omega_o$ put $f=p=0$, $s=N_o$, $W=1$, and $\bar g=g_o$;
inside put $\bar g=g+s^2\dd f^2$. On either side set
$x=sw$, $h=x\bar g$, and $A=Ns\bar g^{-1}$, with $N=N_i$ or $N_o$.
Here $d$ is a scalar, whereas $\dd$ denotes differentiation.

For a finite height $M\ge1$ and a homotopy parameter $u_*\in[0,1]$ the
equations are
\begin{equation}\label{eq:compact-system}
\begin{aligned}
 \divg_g p&=TN_i &&\text{on }\Omega_i,\\
 B&=\nabla N_i-(u_*K+\ell g)p,\qquad
 \ell=T-u_*\tau &&\text{on }\Omega_i,\\
 B&=\nabla N_o,\qquad \Sigma_-=0 &&\text{on }\Omega_o,\\
 \divg B&=\Sigma_++\Sigma_-,\qquad
 -\divg(A\dd w)=\Sigma_+x+P &&\text{on each side}.
\end{aligned}
\end{equation}
We first truncate the exterior at a smooth distant coordinate sphere
$\mathcal S_R$. The complete boundary conditions are
\begin{equation}\label{eq:compact-boundary}
\begin{array}{c|l}
 S&f=u_*M,\quad B_n=-b_*,\quad (A\dd w)_n=0,\\
 \mathcal C&f=0,\quad N_o=s_i,\quad (B_i)_n=(B_o)_n,
       \quad w_i=w_o,\quad (A_i\dd w_i)_n=(A_o\dd w_o)_n,\\
 \mathcal S_R&N_o=w=1.
\end{array}
\end{equation}
Here $b_*>0$ is fixed throughout the homotopy. Eventually $R\to\infty$
with every other parameter fixed, and the last row becomes $N_o,w\to1$.

We specify all cutoffs used in these equations. Let $d_S$ be a fixed
smooth function, positive away from $S$ and equal to inward distance on
a collar of $S$. Set $F=u_*M-f$. The quotient $F/d_S$ at $S$ is its
continuous extension. For example, in that collar it equals
$\int_0^1\partial_nF(q,td_S)\,\dd t$; hence it has one fewer controlled
derivative than $F$ in the relevant H\"older spaces. Choose smooth
cutoffs whose product $\chi_T$ is supported where
\begin{equation}\label{eq:compact-trace-cutoff}
 F<M/10,\qquad z<1/2,\qquad FN_i\sqrt{x}<M^{-1},\qquad
 N_iF/d_S<M^{-1},
\end{equation}
and equals one if all four upper bounds hold with an additional factor
$1/2$. The slope cutoff is a smooth function of $z^2$. The other scalar
cutoffs are extended as one for sufficiently negative arguments. Define
\begin{equation}\label{eq:compact-sources}
 T=u_*\tau(1-\chi_T),\qquad
 \Sigma_-=-\tfrac12N_i\ell^2\zeta(N_i),\qquad
 \Sigma_+=u_*LW\eta_M\chi_x(x)\chi_N(N/\eps).
\end{equation}
The fixed smooth function $\zeta$ lies in $[0,1]$, equals one for
$N\le\delta_0$, and vanishes for $N\ge2\delta_0$.
Here $0<\delta_0,b_*,\delta<1$, $L\ge1$, $0<\eps<1$;
$\chi_x=1$ on $x\le1-\delta$ and is zero on $x\ge1-\delta/2$,
whereas $\chi_N=0$ on $N\le\eps$ and is one on $N\ge2\eps$.
In particular
\begin{equation}\label{eq:compact-source-size}
 u_*\tau\le T\le0,\quad |\ell|\le C,\quad
 -C N_i\le\Sigma_-\le0,\quad
 \|\Sigma_-\|_{L^1}\le C\delta_0,\quad 0\le\Sigma_+\le LW.
\end{equation}

The spatial cutoff $\eta_M$ is chosen as follows. Let $\psi_\infty$
be the exterior capacitary potential, equal to one on $\mathcal C$ and
zero at infinity. Fix a large radius $R_\delta$ beyond which
$\psi_\infty<\delta$. Such a potential and radius follow from scalar
Dirichlet theory and the barriers $C|x'|^{-\beta}$, $0<\beta<1$,
on the fixed asymptotically flat end. Fix an exterior collar width
$a_o>0$. The designated set $\mathcal D_M$ consists of all of
$\Omega_i$ except an inner seam collar of width $2M^{-4}$, together
with the exterior between collar distance $2a_o$ and radius
$R_\delta$. Enlarge $R_\delta$ if necessary to contain that collar.
Take $\eta_M=1$ on a neighborhood of $\mathcal D_M$, supported away
from the seam, with inner separation at least $M^{-4}$ and exterior
separation at least $a_o$, and with support in a common compact region.
All spatial cutoffs below have slightly nested supports of this kind.
Thus $N_o$ is harmonic on the fixed exterior seam collar. The choice of
$R_\delta$ does not involve an unknown solution: the exterior equations
give $N_o\Delta_{g_o}x=-P\le0$, so comparison gives
\begin{equation}\label{eq:compact-capacity-comparison}
 x\ge1-\psi_\infty \quad\text{on }\Omega_o.
\end{equation}
The same comparison holds on truncations since the right side is at
most one on $\mathcal S_R$. Consequently every exterior point with
$x<1-\delta$ lies inside the fixed radius $R_\delta$.

There are two further sources, $P=u_*(P_1+P_2)$. Put
$T_1=M^{9/10}$ and $T_2=C_2/\eps$. The source $P_1$ is an adjustable
finite constant $D_1$ times smooth cutoffs, supported inside where
$f>M/2$ and $w<4T_1$, and equal to $D_1$ where
$f\ge4M/5$ and $w\le2T_1$. The source $P_2$ is $D_2$ times smooth
cutoffs on the two spatial supports just described, vanishes if
$N\ge4\eps$ or $w\ge4T_2$, and equals $D_2$ on a neighborhood of
$\mathcal D_M$ where $N\le3\eps$ and $w\le2T_2$.
The constants $C_2,D_1,D_2$ will be chosen below. Every equation is a
smooth finite equation on the open domain $N_i,N_o,w>0$.

\begin{proposition}[Finite compact construction]\label{prop:compact-solutions}
Fix base data satisfying \eqref{eq:compact-base-data}. Fix
$b_*,\delta_0,\delta,a_o>0$ and $L\ge1$. For every sufficiently large
finite $M$, $C_2$ can be chosen depending on these parameters, and then
$\eps>0$ can be chosen arbitrarily small and $D_1,D_2$ sufficiently
large but finite, so that \eqref{eq:compact-system}--\eqref{eq:compact-boundary}
at $u_*=1$ has a positive solution on the full attached exterior.
It is smooth on each closed side of $\mathcal C$ and tends to the
specified values at infinity. It satisfies
\begin{equation}\label{eq:compact-bounds}
 \begin{gathered}
 0\le f\le M,\qquad N_i,N_o\le C(L)M,\qquad
 \int_{\Omega_i}W^2\le C(L)M^2,\\
 N_i\ge c\eps,\qquad s_i,N_o\ge c(M,L)\eps.
 \end{gathered}
\end{equation}
The constants in these bounds do not depend on $\eps,D_1,D_2$ or the
outer truncation. Dependence on the fixed data and on the preceding
parameters is suppressed. At fixed $M,L,\eps$, the lapse and height
have uniform local continuity bounds off the seam, including at $S$,
independent of $D_1,D_2$. Moreover
\begin{equation}\label{eq:compact-penalty-targets}
 w\ge M^{9/10}\quad\text{where }f\ge9M/10,
 \qquad x\ge1\quad\text{on }\mathcal D_M\cap\{N\le2\eps\},
\end{equation}
and
\begin{equation}\label{eq:compact-penalty-cost}
 \int P_1\le C(L)M^{-1/10},\qquad
 \int P_2\le\omega_M(\eps),\qquad \omega_M(\eps)\longrightarrow0
 \quad(\eps\downarrow0).
\end{equation}
The function $\omega_M$ is uniform in the finite strengths and outer
radius. In particular
\begin{equation}\label{eq:compact-bad-volume}
 \int_{\mathcal D_M\cap\{x<1-\delta\}}W\le C/L,
\end{equation}
where $C$ is independent of $L,M,\eps,D_1,D_2$. On the smooth sides the
metric $h=x\bar g$ obeys the scalar and boundary identities
\eqref{eq:compact-scalar} and \eqref{eq:compact-mean-curvature} below.
Its asymptotically flat mass satisfies
\begin{equation}\label{eq:compact-mass-budget}
 m_h\le m_o+C(b_*+\delta_0)+C(L)M^{-1/10}+C\omega_M(\eps).
\end{equation}
Thus the parameter order is: fix the base data; choose $b_*,\delta_0$
for the mass tolerance; choose $\delta,a_o$ and then $L$ for the area
tolerance; choose $M$ large; choose $C_2$; choose $\eps$ small; choose
the finite strengths; and finally let the outer radius tend to infinity.
\end{proposition}

We prove estimates for positive zeros in the H\"older domain used for
continuation in Section~\ref{sec:transmission-analysis}, before invoking
existence.  Thus $f\in C^{3,\alpha'}(\overline\Omega_i)$ and
$N,w\in C^{2,\alpha'}$ on each closed side, for the fixed positive
exponent $\alpha'$ selected there.  This is membership in the domain,
not an a priori bound on its norm.  No a priori derivative bound on a
cutoff is used in the following finite estimates.

\begin{lemma}[Lapse and slope bounds]\label{lem:compact-lapse}
Every such positive zero, for $0\le u_*\le1$ and sufficiently large
outer radius, satisfies \eqref{eq:compact-bounds}, with $M$ replacing
the upper height bound by $u_*M$. These estimates are independent of
the strengths of $P_1,P_2$ and of all derivatives of $T$.
\end{lemma}

\begin{proof}
The minimum principle for the height equation gives $f\ge0$ because
$T\le0$. A maximum exceeding $u_*M$ has $F<0$ and $z=0$, so every
trace cutoff is fully active in its neighborhood and $T=0$ there.
The strong maximum principle, applied on the exceeding set, excludes
such a maximum. Hence $0\le f\le u_*M$.

Let $\psi_R$ be harmonic on the truncated exterior, one on
$\mathcal C$ and zero on $\mathcal S_R$, and set
$\gamma_R=-\partial_n\psi_R|_{\mathcal C}$. The functions $\gamma_R$
have uniform positive upper and lower bounds. Also $\psi_R$ has a
uniform positive lower bound on all exterior source supports. To check
uniformity, compare $\psi_R$ from below with a fixed finite annular
capacitary solution and from above with $\psi_\infty$, and apply
boundary estimates and the Hopf lemma on a fixed collar. Green's
formula on the exterior and the divergence theorem on the inside give
\begin{equation}\label{eq:compact-capacity}
 \int_{\Omega_i}\Sigma_+
 +\int_{\Omega_o}\psi_R\Sigma_+
 +\int_{\mathcal C}\gamma_Rs_i
 =\int_{\mathcal C}\gamma_R+\int_S b_*
      -\int_{\Omega_i}\Sigma_-\le C.
\end{equation}
For example the seam contribution in the exterior Green formula is
$-B_n-\gamma_RN_o$, and its far contribution is
$-\partial_n\psi_R$; the harmonic flux identity for $\psi_R$ gives
the displayed signs. The constant uses only
$\int|\Sigma_-|\le C\delta_0$. In particular $\min_{\mathcal C}s_i\le C$.

Testing the height equation with $f$ and using its two constant traces
gives, since $p\cdot\dd f=W^2-1$,
\begin{equation}\label{eq:compact-height-energy}
 \int_{\Omega_i}(W^2-1)
 =-u_*M\int_S p_n-\int_{\Omega_i}fTN_i
 \le CM\sup_{\Omega_i}N_i.
\end{equation}

For the lower lapse estimate freeze
$Z=(u_*K+\ell g)v$, which satisfies $|Z|\le C$, and freeze the positive
seam function $d=s_i/N_i\le1$. Solve the homogeneous density problem
with current $\nabla\widehat N_i-Z\widehat N_i$ inside, ordinary
gradient outside, divergence zero, zero current flux at $S$, matching
current at the seam, $\widehat N_o=d\widehat N_i$, and outer value one.
This is an independent linear problem. Its scalar Fredholm assertion
is proved in Lemma~\ref{lem:transmission-fredholm}, using only frozen
linear estimates, without the present nonlinear estimates. Its adjoint
has operator $\Delta+Z\cdot\nabla$ inside, ordinary Laplacian outside,
ordinary Neumann condition at $S$, continuous value at $\mathcal C$,
and $\partial_n h_i=d\partial_n h_o$ there. The maximum and Hopf
principles give zero adjoint kernel, hence invertibility. Dual testing
with nonpositive scalar forcing gives $\widehat N\ge0$. Harnack and
the Hopf lemma give strict positivity, also at the seam: a zero there
would be a zero on both sides and the opposite Hopf derivative signs
contradict current matching. The H\"older domain gives
$d\in C^{2,\alpha'}$ and $Z\in C^{1,\alpha'}$ (in fact
$Z\in C^{2,\alpha'}$), so the independent scalar lemma applies
separately to each frozen problem. Its inverse may depend on these
coefficient norms. The uniform estimate below uses only $|Z|\le C$
and $0<d\le1$ and does not use those inverse norms.

Put $H_*:=\sup\widehat N_i$. The homogeneous version of
\eqref{eq:compact-capacity} implies $H_*\ge1$, and exterior comparison
gives $\widehat N_o\le H_*$. We claim, uniformly in the frozen fields,
\begin{equation}\label{eq:compact-homogeneous-harnack}
 \inf_{\Omega_i}\widehat N_i\ge cH_*.
\end{equation}
Here it is essential that no lower bound on $d$ is required. In Fermi
coordinates on the two sides reflect the normalized exterior field
into the inside and add it to $\widehat N_i/H_*$. The two coefficient
matrices, including volume densities, agree on the face and differ by
$O(|t|)$ at distance $|t|$ from it. For the positive exterior harmonic
field, interior gradient estimates and Harnack on balls of radius
comparable to $|t|$ give
$|t|\,|\nabla\widehat N_o|\le C\widehat N_o$.
Consequently the sum $U$ satisfies a uniformly elliptic divergence
equation with flux correction bounded by $CU$ and zero natural flux
on the face. Expressing that correction as a bounded vector times $U$
and reflecting again gives the ordinary divergence drift Harnack
inequality. Harnack chains on the fixed inner region, together with
this collar estimate and normalization by the inner supremum, give a
uniform positive lower bound for $U$ on the seam. The inner trace is
at least $U/2$, because $\widehat N_o=d\widehat N_i\le\widehat N_i$.

To carry this trace bound uniformly into a collar, take a smooth small
half-domain of radius $h$ in an inner chart and replace the normalized
inner field by its harmonic function with the same boundary data.
The replacement error is at most $Ch$: the divergence forcing
$Z\widehat N_i/H_*$ is bounded, and scaled zero-Dirichlet
$W^{1,p}$ estimates with $p>3$ give precisely that bound. The harmonic
replacement has a fixed positive lower bound centrally up to its
flat face, by the lower face data and nonnegative remaining data.
First choose $h$ small enough and then use interior Harnack chains.
This proves \eqref{eq:compact-homogeneous-harnack}.

The ratio $r_N=N/\widehat N$ is continuous across the seam. On a
sublevel $r_N<\eps/H_*$ we have $N<\eps$, so $\Sigma_+=0$ and
$\Sigma_-\le0$. Dividing the density equation by $\widehat N$ gives
a scalar supersolution inequality with no zeroth order term.
At the seam its normal derivatives match with positive weights
$\widehat N_i,\widehat N_o$; the flux $-b_*$ at $S$ has the favorable
sign for a minimum. The minimum and boundary point principles exclude
a minimum below $\eps/H_*$. Therefore
$N_i\ge(\eps/H_*)\widehat N_i\ge c\eps$.

Next suppose, with $L$ fixed, that the upper estimate fails. Along a
sequence write $N_*:=\sup N_i$ with $N_*/M\to\infty$. The exterior
maximum is bounded by $\max(N_*,1)$ since $\Delta N_o\ge0$.
By \eqref{eq:compact-height-energy},
\[
 \|\Sigma_+/N_*\|_{L^2(\Omega_i)}
 \le CL\bigl(N_*^{-2}+M/N_*\bigr)^{1/2}\longrightarrow0.
\]
The normalized negative source also tends to zero, and the exterior
source tends uniformly to zero. Equation~\eqref{eq:compact-capacity}
shows that $N_o/N_*$ has seam trace tending to zero in $L^1$.
The normalized exterior field tends to zero on every open exterior
compact set. Indeed it is bounded above by $N_*^{-1}$ plus the positive
harmonic extension, decaying at infinity, of its nonnegative seam
trace. That extension is bounded by $\psi_\infty$ and tends locally
to zero off the seam by the $L^1$ trace convergence, using a fixed
annular Poisson estimate and then the decay of $\psi_\infty$.

The reflected sum of the actual normalized lapses is uniformly
H\"older continuous up to the seam by the preceding reflected-sum
argument, now with scalar forcing tending to zero in $L^2$.
Interior Laplace divergence estimates give the same compactness away
from the seam and at $S$. After passage to a subsequence the open
inner limit $V$ extends continuously to the seam using the limit of
the sum. It is nonnegative and nonzero: an inner supremum equal to one
either lies away from the seam or gives a sum at least one in the
seam collar. With $Z$ converging weakly star, this limit satisfies the
homogeneous density equation with natural zero flux on both inner
faces. Here is a way to justify the latter assertion without individual
inner trace compactness. Use inner smooth test functions with ordinary
zero normal derivative at both faces, extended across the seam with
matching values and outer support in the harmonic collar. Integrate
the lapse derivatives onto the tests. All outer volume and trace
terms vanish by the convergence just proved, as do the normalized
sources and $b_*$. Thus
\[
 \int_{\Omega_i}\bigl(V\Delta\varphi+VZ_\infty\cdot\nabla\varphi\bigr)=0
 \quad\text{for every smooth ordinary Neumann test }\varphi.
\]
Solve the Neumann Laplace problem with divergence datum $VZ_\infty$.
Testing the difference by Neumann Poisson solutions shows that its
solution differs from $V$ by a constant. Hence $V\in H^1$ and the
natural zero-flux equation holds in the usual weak sense. Reflection
and Harnack show $V>0$ on the whole closed inner region. In particular
the original normalized inner fields are bounded below away from the
seam; near it the sum, the trace inequality, and the harmonic
replacement argument above give the same lower bound. The replacement
error now has the additional vanishing scaled $L^2$ forcing term.

We reach a contradiction by the differential estimate of
Lemma~\ref{lem:gradient-maximum}. Put $G=-\log s$ and maximize
$G+j(f)$ on $\Omega_i$, where
$j(f)=D(f^2-2Mf)$ and $D>0$ is fixed sufficiently large depending on
$L$. Since $j\le0$ on the height range and $j(0)=0$,
at the maximum $s\le\min_{\mathcal C}s\le C$. Since
$N_i\ge cN_*$ there, $y\to1$. At an interior maximum that lemma gives
\begin{equation}\label{eq:compact-gradient-test}
 j''q|\dd f|^2\le C\{1+\Sigma_+/N_i+
                    (1-y)^2(j'|\dd f|)^2\},\qquad
 q=\frac{1-y}{1+y}.
\end{equation}
Use $q|\dd f|^2=J/s^2$, $J=y/(1+y)$,
$\Sigma_+/N_i\le L/s$, and
$(1-y)^2(j'|\dd f|)^2\le(j'/N_i)^2=o(1)$.
Multiplication by $s^2$ contradicts \eqref{eq:compact-gradient-test}
once $D$ exceeds a fixed multiple of $1+L$ and the fixed upper bounds
for $s^2$ and $s$ at this maximum.

At a constant-height boundary, write $\widehat z=v\cdot n\le0$.
Lemma~\ref{lem:graph-boundary} gives
\begin{equation}\label{eq:compact-boundary-gradient}
 d^2\partial_n(G+j(f))
 =-B_n/N_i+u_*k_T\widehat z-yH+j'\widehat z/N_i.
\end{equation}
At $S$ this is strictly positive because $j'\le0$, $H\le0$,
$k_T<0$, and $B_n=-b_*$. That is impossible for an inward derivative
at a maximum. At a seam maximum $s$ is its seam minimum. The positive
harmonic exterior lapse on the fixed collar satisfies
$B_n\ge-Cs$ at that point: compare with the seam minimum times a
fixed harmonic function equal to one on the seam and zero on the
other collar boundary. Equation~\eqref{eq:compact-boundary-gradient}
is then at most $Cd+|k_T|z-yH_i+o(1)<0$, using
$H_i>|k_T|$. That is impossible for an outward derivative at a maximum.
This proves $N_i,N_o\le C(L)M$ and, by
\eqref{eq:compact-height-energy}, the stated $L^2$ estimate for $W$.

Finally keep $M,L$ fixed and suppose $\inf s_i/\eps\to0$.
Choose a fixed smooth $j$ on $[0,M]$ with $j''>0$ and
$0<j'<b_*/2$. At a maximum of $G+j(f)$, bounded oscillation of $j$
implies $s_i/\eps\to0$. The lapse floor gives $y\to1$.
Multiply \eqref{eq:compact-gradient-test} by $s^2$: its left side
stays positive, while its right side is bounded by
$C(s^2+Ls+(j')^2d^2)$ and tends to zero. At $S$, the first and last
terms of \eqref{eq:compact-boundary-gradient} have sum at least
$(b_*-j')/N_i>0$. At the seam its last term is nonpositive, and the
same harmonic collar comparison makes the derivative strictly
negative. These contradictions prove $s_i\ge c(M,L)\eps$.
On the exterior, below the level $\eps$ the positive source vanishes
and the lapse is harmonic. Its boundary values and the minimum
principle therefore give $N_o\ge c(M,L)\eps$ too.
\end{proof}

\begin{lemma}[Finite penalty selection]\label{lem:compact-penalties}
For fixed preceding parameters the sources can be selected with finite
strengths to obtain \eqref{eq:compact-penalty-targets}--\eqref{eq:compact-bad-volume}
in every actual positive solution, uniformly in the outer truncation.
Their costs obey \eqref{eq:compact-penalty-cost}.
\end{lemma}

\begin{proof}
All assertions in this proof concern $u_*=1$. First choose
$C_2\ge c(M,L)^{-1}$ from Lemma~\ref{lem:compact-lapse}, so that
$sT_2\ge1$. This precedes the choice of $\eps$. For a piecewise
Lipschitz test $\varphi$ with common seam value and zero outer value,
testing the positive $w$ equation with $\varphi^2/w$ gives
\begin{equation}\label{eq:compact-penalty-test}
 \int P\frac{\varphi^2}{w}\le\int A(\dd\varphi,\dd\varphi).
\end{equation}
Indeed the omitted term $\Sigma_+x\varphi^2/w$ is nonnegative,
and completing the square in the two derivative terms gives the
right side. Flux matching cancels the seam terms and the inner flux
is zero.

Take $\varphi=f/M$, extended by zero outside. Directly from
\eqref{eq:compact-constitutive},
$A(\dd f,\dd f)=W-W^{-1}$. On the support of $P_1$ we have
$\varphi>1/2$ and $w<4T_1$. Thus
\[
 \int P_1\le\frac{16T_1}{M^2}\int_{\Omega_i}W
 \le C(L)M^{-1/10},
\]
by Cauchy--Schwarz and \eqref{eq:compact-bounds}.

For $P_2$ take spatial localizers equal to one on its supports and
vanishing on slightly larger collars of the seam, multiplied by a
scalar cutoff of $N/\eps$ which is one below four and zero above five.
These tests may meet $S$ and have zero seam value on each side. Since
$A\le N^2g_{\rm base}^{-1}$ and $w<4T_2$ on the source support,
\begin{equation}\label{eq:compact-low-lapse-cost}
 \int P_2\le C_M\eps+
 C_M\eps^{-1}\int_{\mathcal U\cap\{N<5\eps\}}|\nabla N|^2.
\end{equation}
The fixed localization sets $\mathcal U$ stay away from the seam;
constants here may depend on $M$ and every preceding parameter.
Testing the lapse equation with a squared spatial cutoff times
$(6\eps-N)^+$ gives
\begin{equation}\label{eq:compact-low-lapse-energy}
 \int_{\mathcal U\cap\{N<5\eps\}}|\nabla N|^2
 \le C_M\eps^2+
 C\eps\int_{\mathcal U'\cap\{N<6\eps\}}\Sigma_+.
\end{equation}
To check the signs, on this sublevel the current correction has size
at most $CN\le C\eps$, so Young's inequality absorbs its gradient
term and the cutoff terms cost $C_M\eps^2$. The negative source can
be discarded. At $S$ the outward current equals $b_*>0$, so its
contribution decreases the right side of the energy inequality.

We next prove that the last integral tends uniformly to zero as
$\eps\downarrow0$. This step does not use a positive lapse floor.
If uniformity failed, take a violating sequence. Equations
\eqref{eq:compact-height-energy} and \eqref{eq:compact-bounds} bound
the positive sources in $L^2$, independently of $\eps$ and the
strengths. The lapse equation with bounded divergence datum $ZN$
therefore gives local $W^{1,p}$ bounds for every $3<p<6$, also at $S$
with the fixed flux data. After passing to a subsequence, the lapses
converge uniformly to a nonnegative function $V$ on the localization
sets, the drifts converge weakly star, and the sources converge weakly
in $L^2$. The limit equation has the form
\[
 \Delta V=\divg(Z_\infty V)+g_1.
\]
The negative-source limit vanishes almost everywhere on $\{V=0\}$
because $|\Sigma_-|\le CN$.

For completeness the scalar source $g_1$ also vanishes there. Let
$q$ be a density point of $\{V=0\}$ and a Lebesgue point of $g_1$
and $|g_1|^2$ in an interior chart. Rescale a ball of radius $r$ to a
unit ball and divide $V$ by $r^2$. The resulting nonnegative functions
vanish at the center. The inhomogeneous Harnack estimate gives a
uniform bound on a smaller ball: the rescaled drift is $rZ_\infty$,
and the rescaled scalar forcing has bounded $L^2$ norm at this
Lebesgue point. Since the zero sets occupy a fraction tending to one,
the rescaled functions tend to zero in local $L^1$. Testing their
equations against a fixed smooth bump, with derivatives transferred
to that bump, gives $g_1(q)=0$. Thus the nonnegative positive-source
limit $\sigma_+$ vanishes almost everywhere on $\{V=0\}$ as well.
Boundary points do not matter for these volume integrals. For each
$t>0$, uniform convergence implies
$\{N<6\eps\}\subset\{V\le t\}$ along the sequence eventually.
Weak $L^2$ convergence, tested against the characteristic function of
this fixed set, and then $t\downarrow0$, show
\[
 \limsup\int_{\mathcal U'\cap\{N<6\eps\}}\Sigma_+
 \le\int_{\mathcal U'\cap\{V=0\}}\sigma_+=0.
\]
Combining this with \eqref{eq:compact-low-lapse-cost} and
\eqref{eq:compact-low-lapse-energy} proves the second cost in
\eqref{eq:compact-penalty-cost}, uniformly in both strengths.

Now fix such an $\eps$. The lapse and slope bounds give uniform
ellipticity, independent of strengths. They also give uniform
continuity of $N$ and $f$ on the localization sets. For $N$, use the
same Laplace divergence estimates, now with bounded scalar forcing;
for $f$ the already established gradient bound suffices. Consequently
each point with $f\ge9M/10$ has a neighborhood of a uniform positive
radius on which $f\ge4M/5$. It is uniformly separated from the seam
unless the neighborhood ends at $S$, where the constant trace permits
reflection. Likewise every point of
$\mathcal D_M\cap\{N\le2\eps\}$ has a uniform neighborhood within
the full spatial cutoff range and with $N\le3\eps$.

In either neighborhood let $T_j$ be its corresponding threshold.
The volume of $\{w\le2T_j\}$ there is at most
$D_j^{-1}\int P_j$, hence tends uniformly to zero as $D_j\to\infty$.
The function $(2T_j-w)^+$ is a bounded weak subsolution of the
homogeneous scalar divergence equation with coefficient $A$.
Its local supremum is bounded by its $L^2$ mean, with a uniform
constant on these fixed neighborhoods; at $S$ use the homogeneous
natural-flux reflection. Since its support has the preceding small
volume and its value is at most $2T_j$, this supremum is strictly
less than $T_j$ once the strength is sufficiently large. This gives
$w\ge T_1$ on the height target and $w\ge T_2$ on the low-lapse
target. The latter implies $x\ge1$ since $sT_2\ge1$.
The choices can be made with a strict margin, so the weak inequalities
survive limits of outer truncations.

If $x<1-\delta$ at a designated point, the second target now forces
$N>2\eps$. All cutoffs in $\Sigma_+$ equal one, so
$\Sigma_+=LW$ there. Apply \eqref{eq:compact-capacity}, using the
positive lower bound of $\psi_R$ on the common exterior support.
This proves \eqref{eq:compact-bad-volume} with a constant independent
of $L,M,\eps$ and the strengths.
\end{proof}

\begin{proof}[Proof of Proposition~\ref{prop:compact-solutions}]
Choose the parameters in the stated order, using
Lemmas~\ref{lem:compact-lapse} and \ref{lem:compact-penalties}.
Proposition~\ref{prop:transmission-existence}, proved independently in
the next section, now supplies a positive zero for each sufficiently
large compact truncation. Its hypotheses are exactly the finite
smooth cutoffs above, the homotopy bounds of
Lemma~\ref{lem:compact-lapse}, and the fixed exterior harmonic collar.
Its continuation starts with the ordinary scalar transmission problem
at $u_*=0$; it does not assume solvability of the coupled problem.
That proposition also supplies an expanding-truncation subsequence
converging smoothly on compact subsets of each closed side, with
$N_o,w\to1$ at infinity. All the bounds and penalty conclusions pass
to this limit. This is the only invocation of nonlinear existence in
the present proof.

We record the geometric conclusions precisely. Set
$\pi=N_i^{-1}\operatorname{sym}\nabla p^\flat$,
$\lambda=\pi-K$ inside, and set these fields to zero outside.
Lemma~\ref{lem:graph-current}, applied at $u_*=1$, gives
\begin{equation}\label{eq:compact-scalar}
 xR_h=D(v)+|\lambda|^2-\ell^2-2\Sigma_-/N
       +2P/(Nx)+\tfrac32|\dd\log x|_{\bar g}^2.
\end{equation}
On the exterior take $D=R_{g_o}=0$. The current
$Q=B+A\dd w$ has matching normal flux and
$\divg Q=(1-x)\Sigma_++\Sigma_- -P$. At a constant-height face,
Lemma~\ref{lem:graph-boundary} gives
\begin{equation}\label{eq:compact-mean-curvature}
 s\sqrt{x}\,H_h
 =N(H-k_T\widehat z)+B_n+(A\dd w)_n/x.
\end{equation}
At $S$ this is strictly negative. At $\mathcal C$ the induced metrics
agree, because $s_i=N_o$ and $w$ is continuous. The seam mean
curvatures satisfy $H_{h,i}\ge H_{h,o}$ in the specified direction:
the flux terms agree and
\[
 H_i-k_T\widehat z\ge
 \sqrt{1-\widehat z^2}\sqrt{H_i^2-k_T^2},
\]
as follows by squaring, since the difference of squares is
$(H_i\widehat z-k_T)^2$ and both relevant sides are nonnegative.

Beyond the fixed source supports, both $N_o$ and $x=N_ow$ are
harmonic and tend to one. Radial barriers and scaled estimates give
$N_o-1,x-1=O_2(|x'|^{-\beta})$ for every fixed $1/2<\beta<1$.
Thus $h=xg_o$ is asymptotically flat, and the flux of $Q$ equals
the limiting flux of $\nabla x$; the remaining terms have vanishing
integral by $2\beta>1$. The linear conformal ADM variation is
$m_h-m_o=-(8\pi)^{-1}\lim\int\partial_nx$. Integration of the
current equation therefore gives the exact identity
\begin{equation}\label{eq:compact-mass}
 8\pi(m_h-m_o)
 =\int_S b_*+\int\bigl[(x-1)\Sigma_++P-\Sigma_-\bigr].
\end{equation}
Since the support of $\Sigma_+$ has $x<1$, discarding its contribution
and using the source costs proves \eqref{eq:compact-mass-budget}.
The scalar curvature decays as $O(|x'|^{-2-2\beta})$ on the end,
by \eqref{eq:compact-scalar}; this has integrable decay since
$2+2\beta>3$.
\end{proof}

The omitted collars have a controlled cost in the area comparison.
Indeed $W=1$ outside, and the inner omitted collar has volume
$O(M^{-4})$. Hence
\begin{equation}\label{eq:compact-full-bad-volume}
 \int_{\{x<1-\delta\}}W
 \le C/L+C a_o+C(L)M^{-1}.
\end{equation}
This uses \eqref{eq:compact-capacity-comparison} outside $R_\delta$
and Cauchy--Schwarz with \eqref{eq:compact-bounds} in the inner
collar. It explains why $a_o$ and $L$ are chosen before $M$.
Finally the possible negative scalar curvature is confined to the
closed set defined by the weak inequalities in
\eqref{eq:compact-trace-cutoff} and $N_i\ge\delta_0$.
Outside that set either $\ell=0$ or $\zeta=1$, and
\eqref{eq:compact-scalar} is nonnegative. On that set the height
penalty and $z\le1/2$ give
$x\ge(\sqrt3/2)\delta_0M^{9/10}$, so
$R_h\ge-C\delta_0^{-1}M^{-9/10}$ everywhere on the smooth sides.
The geometric treatment of this exceptional set is the content of
Section~\ref{sec:scalar-error}.
